\section{Models, equilibria and labels}
\label{sec:methods}

\subsection{Model ensembles}
\label{sec:ensembles}

\begin{table*}
\caption{Model ensembles. ESN: echo state network; mlp, mlp2: perceptrons with one and two hidden layers; res: residual
network; NG-RC: next-generation reservoir computer; ``all'': the seven regimes of Table~\ref{tab:systems} (pooled statistics use six, with LE1 reported separately). H0: noise-free H1. Training sets
are defined in Sec.~\ref{sec:data}; ``light'' and ``full'' refer to the labeling protocols of Sec.~\ref{sec:protocol}; $W^u$: unstable manifolds traced
(Fig.~\ref{fig:schematic}).}
\label{tab:ensembles}
\begin{ruledtabular}
\begin{tabular}{llllll}
Ensemble & Models & Systems & Data, size & Labels & Sec. \\
\hline
Size sweep & $700$ ESN & all & H1, HE1; $N=50$--$2000$ & light & \ref{sec:esn_law} \\
Full-protocol ensemble & $1138$ ESN & all & H1, HE1, E1, H0; $N=300$--$1000$ & full & \ref{sec:theory_checks}, \ref{sec:esn_law}, \ref{sec:e1} \\
Settings sweep & $200$ ESN & four & HE1; $N=300$ & light & \ref{sec:esn_law} \\
Intervention & $82$ ESN & all & H1 $+$ $0$--$20$ transients; $N=300$, $700$ & light & \ref{sec:data_effects} \\
Benchmark & $754$ hard-tanh ESN & Lorenz & one trajectory; $N=8$--$20$ & exact, full & \ref{sec:theory_checks} \\
No-equilibrium control & $32$ ESN & Li--Sprott & one trajectory; $N=300$ & full & \ref{sec:theory_checks} \\
Model classes & $10\,080$ mlp, mlp2, res & all & H1, HE1, E1; $H=8$--$1024$ & light & \ref{sec:theory_checks}, \ref{sec:classes}, \ref{sec:e1} \\
Example & $20$ mlp & Chua~I & HE1; $H=32$ & $W^u$ & \ref{sec:manifold} \\
Penalty & $840$ mlp2 ($720$ analyzed) & all & H1, HE1; $H=256$, $\lambda_{\mathrm{P}}=0$, $0.1$, $1$ & light & \ref{sec:penalty} \\
E1 controls & $210$ NG-RC, $210$ quadratic & all & subsets of E1 & full, light & \ref{sec:e1} \\
Positive controls & $180$ NG-RC, $360$ quadratic & six & windows of H1, HE1 & full, light & \ref{sec:e1} \\
Saturation & $184$ ESN, rebuilt & six & H1, HE1; $N=100$--$2000$ & light & \ref{sec:theory_checks} \\
Certified enumeration & $960$ mlp, mlp2, res & six & H1, HE1; $H=64$, $256$ & -- & \ref{sec:theory_checks} \\
\end{tabular}
\end{ruledtabular}
\end{table*}


Table~\ref{tab:ensembles} lists the ensembles. Echo state networks are trained by ridge regression of the readout on the
reservoir states driven by the training series,\cite{LukoseviciusJaeger2009Reservoir} with tanh activation, a sparse random
$W$ of spectral radius $\rho$ and the input scale and density given in Appendix~\ref{app:training}. The \emph{size sweep}
varies $N$ from $50$ to $2000$; the \emph{full-protocol ensemble} is labeled with the full protocol below and also supplies
the indices of Sec.~\ref{sec:theory_checks} and the networks trained on E1; the \emph{settings sweep} varies the input scale
and the reservoir density; the \emph{intervention ensemble} keeps each reservoir and refits only the readout on the data
variants of Sec.~\ref{sec:data}. The \emph{benchmark} consists of hard-tanh networks whose fixed points are enumerated
exactly (Sec.~\ref{sec:equilibria}), and the \emph{equilibrium-free control} repeats the setting of R\"ohm
\textit{et al.}\cite{Rohm2021} on the Li--Sprott system, which has no equilibria. The perceptrons and residual networks
are trained by Adam on a one-step loss followed by a four-step rollout loss, for eight hidden widths from $8$ to $1024$ and
$20$ seeds per configuration, on H1, HE1 and E1. Two further ensembles serve as interventions and controls: perceptrons
trained with a penalty on unstable equilibria (Sec.~\ref{sec:penalty}) and models whose function class contains the Euler
step of the true vector field for four of the systems (Sec.~\ref{sec:e1}).

\subsection{Equilibria and their indices}
\label{sec:equilibria}

Equilibria are found by Newton's method from many starting points: uniform points in the range of the map (or in an
enlarged region for unbounded models), states visited during training, and the images of the true equilibria in the
model's state space; for echo state networks deflation\cite{Farrell2015Deflation} pushes Newton away from roots already found.
Each root is refined to machine precision, and its Jacobian gives the index~\eqref{eq:index}, the number of unstable
directions and the number of real eigenvalues greater than $1$. Root finding in high dimension is never guaranteed to be
complete;\cite{SussilloBarak2013,AllgowerGeorg2003Continuation} Corollary~\ref{cor:counts}(iii) provides the check we use,
and Sec.~\ref{sec:theory_checks} adds two more: searches with up to eight times as many starts and, for the
three-dimensional models, a certified enumeration by interval arithmetic.
The survival analyses of echo state networks use only networks whose found equilibria have index sum $1$; for the
perceptrons we report the effect of the same restriction. For unbounded models the
check does not apply, and points to which Newton's method runs off to infinity ($|x|>10^3$ in physical units, where the tanh
units saturate) are discarded.

To test the root finding and the labels against ground truth, we use echo state networks with hard-tanh activation
$\phi(z)=\max(-1,\min(1,z))$, trained by the same ridge regression with $N=8$--$20$. In each of the $3^N$ regions where
every neuron is saturated below, linear or saturated above, the closed-loop map is affine, so all fixed points follow from
small linear systems and a sign check. Exhaustive enumeration over the regions gives the complete set of fixed points,
which piecewise-linear search heuristics do not guarantee.\cite{Eisenmann2023Bifurcations} We keep the $754$ networks whose
enumerated fixed points have residuals below $10^{-9}$.

\subsection{Reproduction and excitation labels}
\label{sec:protocol}

For each hidden attractor of the true system, the model is started from $64$ states on that attractor; an echo state
network is first synchronized on $100$ samples of the true trajectory and then run in closed loop. Trajectories are
iterated for $2\times10^4$ steps; those that do not settle on an equilibrium, do not escape and have the right type (largest
Lyapunov exponent above or below $0.02$ per unit time) define the model attractor, which counts as reproduced if
$\Delta_{\mathrm{att}}<0.2$ and $d_H<0.2$ (Appendix~\ref{app:protocol}).

The excitation test starts trajectories near every unstable equilibrium of the model and records which of them reach the
reproduced attractor, i.e., end close to an orbit on it after the same number of steps (Appendix~\ref{app:protocol}). Following
Corollary~\ref{cor:criterion}, the points are placed along the unstable eigendirections (both signs, real and imaginary parts
of complex eigenvectors) at radii $10^{-6}$, $10^{-4}$ and $10^{-3}$, together with $50$--$64$ random points of the ball of the
same radius; radii are in the normalized coordinates of the model (reservoir coordinates for echo state networks, physical
coordinates divided by their standard deviation otherwise). The smallest radius at which some point reaches the attractor is
the excitation radius $\hat r$. The attractor is labeled \emph{self-excited} if $\hat r\le10^{-6}$, \emph{excited at
radius $\hat r$} if $10^{-6}<\hat r\le10^{-3}$, and \emph{not excited} otherwise. ``Not excited'' marks a candidate hidden
attractor.

The full protocol (Appendix~\ref{app:protocol}) adds a radius grid from $10^{-6}$ to $10^{-1}$ with $200$ points per radius
and up to $20$ unstable directions, and a search for further equilibria near the basin boundary, where any exciting
equilibrium lies (Corollary~\ref{cor:criterion}), with an upper bound on the probability that this search misses an
equilibrium.\cite{Good1953,McAllesterOrtiz2003,BerendKontorovich2013,BenHamou2017MissingMass} An attractor is labeled
\emph{statistically hidden} when it is not excited at any radius and the bound is at most $0.05$; this notion, defined in
Appendix~\ref{app:protocol}, is how hiddenness can be asserted in a high-dimensional model whose equilibria cannot all be
guaranteed. On the $115$ pairs (network, attractor) of the six regimes reproduced in both protocols the labels agree in all
$115$: the $45$ ``not excited'' attractors are statistically hidden, and every statistically hidden attractor is ``not
excited''; in LE1 they agree in $14$ of $16$. The
light protocol is used for the large ensembles.

\subsection{Statistics}
\label{sec:stats}

The unit of analysis is a pair (model, hidden attractor) with the attractor reproduced; Chua's circuit contributes two
hidden attractors per model. Two analyses use one row per model: the full-protocol survival curve (the best-matching hidden
attractor) and the intervention regression (the first hidden attractor). Shares are reported with Wilson $95\%$
intervals. The dependence of survival on $U$ is modeled by logistic regression of ``not excited'' on $\log(1+U)$ with fixed
effects for the regime and, where stated, terms for the training set, the model class, the hidden width and the network
size;\cite{Agresti2013Categorical,McCullaghNelder1989GLM} the contribution of a group of terms is tested by the likelihood
ratio. Pairs from the same training configuration (regime, set, class, width or size, which share data and differ only in
the seed) are not independent, so we also report cluster-robust standard errors with these configurations as clusters. All
counts and fits are produced by the scripts listed in the Data Availability statement.
